\chapter{SynapseML and AI Integration}
\index{SynapseML}\index{Azure OpenAI}\index{Cognitive Services}\index{Copilot in Fabric}\index{sentiment analysis}\index{batch enrichment}\index{API key management}%

\begin{objectives}
\begin{itemize}
    \item Explain what SynapseML adds on top of Spark MLlib and when to reach for it.
    \item Integrate pre-built AI services---Azure OpenAI and Cognitive Services---into Spark pipelines in batch.
    \item Keep API keys and service credentials out of notebooks through proper secret management.
    \item Enrich a column of customer reviews with sentiment from Azure OpenAI and publish results to Gold.
    \item Use Copilot in Fabric to generate and explain code responsibly.
    \item Design a scalable pipeline that processes terabytes of reviews into BI-consumable sentiment metadata.
\end{itemize}
\end{objectives}

\begin{crosscloud}
Calling managed AI services from distributed data pipelines---enrichment at scale---is a cross-cloud pattern with the same shape everywhere: batch the calls, vault the keys, land the results in the analytical store.

\vspace{2mm}
{\footnotesize
\begin{tabular}{@{}p{2.8cm}p{3.2cm}p{3.2cm}p{3.2cm}p{3.2cm}@{}}
\toprule
\textbf{Capability} & \textbf{Fabric} & \textbf{AWS} & \textbf{GCP} & \textbf{Snowflake} \\
\midrule
Distributed ML library & SynapseML on Spark & SageMaker SDK / EMR MLlib & Vertex AI / Dataproc MLlib & Snowpark ML \\
Managed LLM calls & Azure OpenAI via SynapseML & Bedrock from Lambda/EMR & Vertex AI Gemini APIs & Cortex / external functions \\
Cognitive AI services & Azure AI Services & Comprehend, Rekognition, Textract & Vision, Natural Language, Document AI & Cortex AI functions \\
Secret management & Key Vault / workspace identity & Secrets Manager & Secret Manager & External access integrations \\
AI-assisted coding & Copilot in Fabric & CodeWhisperer/Q Developer & Gemini Code Assist & Copilot integrations \\
Cost model of enrichment & Service calls + CUs & Per-request pricing & Per-request pricing & Credits + external calls \\
\bottomrule
\end{tabular}}

\vspace{1mm}
Databricks runs the same SynapseML library---it originated in the Synapse ecosystem---and Mosaic AI for LLM work; MongoDB Atlas Vector Search occupies the adjacent retrieval niche. The portable rules: \emph{batch instead of row-by-row, vault instead of hardcode, and land enrichment in Gold where BI can reach it}.
\end{crosscloud}

\begin{keyterms}
\textbf{SynapseML}: Microsoft's open-source distributed ML library extending Spark with AI service integrations and additional learners. \quad
\textbf{Azure OpenAI}: managed access to OpenAI models inside Azure compliance and networking boundaries. \quad
\textbf{Batch enrichment}: applying an AI service to a DataFrame column distributed across the cluster, rather than row-by-row. \quad
\textbf{Copilot in Fabric}: the natural-language assistant that generates and explains code in notebooks and other items. \quad
\textbf{Sentiment enrichment}: appending model-computed sentiment columns to review text at pipeline scale.
\end{keyterms}

\section{Week 19 Roadmap}
Week~19 adds managed AI to the data engineer's toolbox. Monday introduces SynapseML. Tuesday covers the pre-built AI services---Azure OpenAI and Cognitive Services---and the secret-management discipline they demand. Wednesday covers Copilot in Fabric as a code-generation assistant, with its review obligations. Thursday is the hands-on project: a SynapseML notebook calling Azure OpenAI to append sentiment to customer reviews. Friday scales the pattern to terabytes of reviews feeding a Gold lakehouse for BI.

The data engineer's role in AI integration is not prompt engineering; it is pipeline engineering around a non-deterministic, rate-limited, per-call-billed dependency \cite{synapseML,microsoftFabricAiServices}. This week teaches that plumbing.

\begin{notebox}
An LLM call inside a Spark pipeline is the most expensive function invocation in your estate by orders of magnitude. Batching, caching, retry policy, and result persistence are not optimizations here---they are the difference between a viable pipeline and a five-figure surprise.
\end{notebox}

\section{Monday: Introduction to SynapseML}
\index{SynapseML!overview}
SynapseML is an open-source library that extends Spark with distributed integrations the base MLlib does not ship: connectors to Azure OpenAI and Cognitive Services, additional learners (for example LightGBM at Spark scale), and utilities for distributed serving \cite{synapseML}. In Fabric it is preinstalled or one environment dependency away, and it runs on the same Spark pools as Chapter~3's notebooks.

\begin{definitionbox}
\textbf{SynapseML} is a Spark-library extension providing distributed transformers for Azure AI services, additional scalable learners, and serving utilities---so AI-service enrichment becomes a DataFrame operation.
\end{definitionbox}

The mental model: SynapseML turns ``call the API'' into a Spark transformer. A column of review text goes in; a column of sentiment comes out; the distribution, batching, and parallelism happen at the executor level. That is precisely what makes it the right tool for terabytes---and what makes its configuration (batch size, concurrency, retry) a production decision rather than a detail.

\section{Tuesday: Azure OpenAI, Cognitive Services, and Secrets}
\index{Azure OpenAI!integration}\index{Cognitive Services}\index{API key management}
\subsection{Pre-Built AI in the Pipeline}
The Azure AI services cover the common enrichment cases: language (sentiment, key phrases, entity recognition), vision, document intelligence, and speech. Azure OpenAI adds generative capability: classification with instructions, summarization, extraction into structured columns \cite{microsoftFabricAiServices,synapseML}. The integration choice is between the pre-built task APIs (cheaper, constrained) and the generative endpoint (flexible, pricier)---start with the task API when it fits.

\subsection{Secret Management}
Every one of these calls authenticates. The rule is absolute: \textbf{no keys in notebooks}. Options in order of preference: Fabric workspace identity / managed identity where the service accepts it; Key Vault references resolved at run time; environment variables injected by the environment or pipeline. Keys never appear in cell text, pipeline JSON, or---critically---in Git, because Week~23 will sync these notebooks to a repository verbatim.

\begin{pitfall}
A key pasted into a notebook ``temporarily'' is a key in Git permanently once Week~23's Git integration syncs the workspace. Treat any key that touched a cell as compromised: rotate it, then fix the plumbing.
\end{pitfall}

\section{Wednesday: Copilot in Fabric}
\index{Copilot in Fabric}
Copilot in Fabric generates and explains code from natural language inside notebooks and other items \cite{microsoftFabricCopilot}. For the data engineer it is a drafting accelerator: boilerplate PySpark, KQL first drafts, unfamiliar API usage, and explanation of inherited code.

The engineering obligations do not change because the author was an assistant. Copilot-generated code gets the same review as any code: Does it read the table it claims to? Does the join key exist? Does it silently drop rows? Is it leakage-safe (Chapter~17's tests still apply)? Does it hardcode a path or secret? Copilot drafts; engineers own.

\begin{tipbox}
Copilot is strongest as a first-draft generator against a precise prompt that names tables, columns, and intent---``read \texttt{silver.reviews}, filter to 2026, dedupe on \texttt{review\_id}, write Delta to \texttt{gold.reviews\_clean}''---and weakest when the prompt is vague enough that you cannot review the output against it.
\end{tipbox}

\section{Thursday: Hands-on Project}
\index{hands-on lab}\index{sentiment analysis!hands-on}
The Week~19 lab uses SynapseML in a Spark notebook to call Azure OpenAI: a column of customer reviews goes in, and a sentiment column comes out, appended and landed for downstream use.

\subsection{Goal and Architecture}
Work in \texttt{fabric-de-week19} with a table \texttt{silver.reviews} (a few thousand synthetic reviews) and an Azure OpenAI deployment (or the instructor-provided endpoint). The notebook batches calls, persists raw responses, and writes the enriched result to \texttt{gold.review\_sentiment}.

\subsection{Step-by-Step Actions}
\begin{enumerate}
    \item Configure the service connection: endpoint and key resolved from Key Vault / environment, never literals.
    \item Build the SynapseML transformer for the chat completion with a constrained prompt: return one of \texttt{positive}, \texttt{neutral}, \texttt{negative}.
    \item Run on a 50-row sample first; inspect raw outputs before scaling.
    \item Set batch size and concurrency deliberately; record the measured cost per thousand rows.
    \item Persist raw responses alongside parsed sentiment (debuggability and reprocessing).
    \item Write the enriched table to Gold with a schema contract: \texttt{review\_id}, \texttt{sentiment}, \texttt{model\_version}, \texttt{scored\_ts}.
    \item Validate: null rate, label distribution sanity, and a hand-checked sample of twenty.
\end{enumerate}

\begin{lstlisting}[language=Python,caption={SynapseML batch sentiment over a reviews DataFrame.},label={lst:chapter19-synapseml}]
import os
from pyspark.sql import functions as F
from synapse.ml.services.openai import OpenAICompletion

reviews = spark.read.table("silver.reviews").where("review_date = '2026-07-31'")

completion = (
    OpenAICompletion()
    .setSubscriptionKey(os.environ["AZURE_OPENAI_KEY"])     # from Key Vault/env
    .setUrl(os.environ["AZURE_OPENAI_ENDPOINT"])
    .setDeploymentName("gpt-4o-mini")
    .setPrompt(
        "Classify the sentiment of this review as exactly one of "
        "positive, neutral, negative. Review: ")
    .setBatchSize(20)                                        # batch, never row-by-row
    .setTextCol("review_text")
    .setOutputCol("raw_response")
)

scored = (completion.transform(reviews.limit(50))
          .withColumn("sentiment", F.lower(F.trim(F.col("raw_response.choices.text")[0])))
          .withColumn("model_version", F.lit("gpt-4o-mini-2026-07"))
          .withColumn("scored_ts", F.current_timestamp()))

(scored.select("review_id", "sentiment", "model_version", "scored_ts")
       .write.format("delta").mode("append")
       .saveAsTable("gold.review_sentiment"))
\end{lstlisting}

\subsection{Validation Checklist}
\begin{itemize}
    \item No key or endpoint appears as a literal anywhere in the notebook.
    \item The 50-row sample run is inspected before any scale-up.
    \item Raw responses are persisted, not just parsed labels.
    \item Batch size and concurrency are set explicitly with measured cost recorded.
    \item The Gold table honors its schema contract including \texttt{model\_version}.
    \item A hand-checked sample agrees with the model at a documented rate.
\end{itemize}

\section{Friday: Use Case and Architecture}
\index{sentiment analysis!at scale}\index{batch enrichment!architecture}
A consumer platform accumulates terabytes of user reviews. Product leadership wants sentiment metadata available in the Gold lakehouse for BI---trends by category, alerts on sentiment drops---processed continuously and affordably.

\subsection{The Pipeline}
New reviews land in Silver through the Part~III machinery. A scheduled Spark job scores the delta (reviews not yet in \texttt{gold.review\_sentiment}, joined on \texttt{review\_id}) with the SynapseML pattern, in bounded batches, within a capacity-friendly window. Results land in Gold with \texttt{model\_version}; the Chapter~8 Direct Lake model picks them up; a Chapter~15 Reflex rule watches the daily negative-share for a sustained climb. Cost control comes from three levers: scoring only the delta, batching aggressively, and routing the bulk to the cheaper task-appropriate deployment.

\begin{figure}[H]
    \centering
    \begin{tikzpicture}[node distance=8mm and 9mm]
        \node[store] (silver) at (0,0) {silver.reviews\\(delta)};
        \node[stage] (syn) at (4.6,0) {SynapseML\\batch scoring};
        \node[store] (gold) at (9.4,0) {gold.review\_\\sentiment};
        \node[stage] (bi) at (13.6,0.9) {Direct Lake\\BI};
        \node[stage, fill=orange!10] (reflex) at (13.6,-0.9) {Reflex\\(sentiment drop)};

        \draw[arrow] (silver) -- (syn);
        \draw[arrow] (syn) -- (gold);
        \draw[arrow] (gold) -- (bi);
        \draw[arrow] (gold) -- (reflex);

        \node[note, text width=12.5cm, align=center] at (7.4,-2.5) {Score the delta, batch the calls, version the model, land in Gold---BI and alerts consume, never recompute.};
    \end{tikzpicture}
    \caption{Terabyte-scale sentiment enrichment feeding BI and alerting.}
    \label{fig:chapter19-sentiment}
\end{figure}

\begin{pitfall}
Re-scoring history on every run is the classic cost explosion: the delta join (``reviews with no sentiment row for the current model version'') is the single most important line in the pipeline. Note the qualifier: a model upgrade legitimately re-scores history, so key the delta on (\texttt{review\_id}, \texttt{model\_version}).
\end{pitfall}

\section{Operational Guidance for Week 19}
Treat AI services as metered external dependencies inside your pipelines: budgeted, batched, cached where legal, retried with backoff, and always persisted with their raw output.

\subsection{Performance and Reliability}
Score deltas, not history. Batch and bound concurrency to the service's rate limits. Retry transient failures with backoff; quarantine persistent failures with the Chapter~10 pattern. Record cost per thousand rows and trend it.

\subsection{Security and Governance Checklist}
\begin{itemize}
    \item Keys resolved from Key Vault or environment; nothing literal, ever.
    \item Review text containing personal data classified before leaving the tenant boundary; the service's data-handling terms reviewed.
    \item Model and deployment versions stamped on every scored row.
    \item Copilot-generated code reviewed with the same rigor as human code.
    \item Rate limits and quotas documented; alerts on throttling.
    \item Raw responses retained per the privacy policy, not indefinitely by default.
\end{itemize}

\section{Interview Q\&A}
\begin{enumerate}
    \item \textbf{Q: What does SynapseML add on top of Spark MLlib?}\\
    A: Distributed transformers for Azure AI services (OpenAI, Cognitive Services), additional scalable learners such as LightGBM, and serving utilities---turning AI-service calls into DataFrame operations.
    \item \textbf{Q: How do you avoid hardcoding Azure OpenAI keys in a notebook?}\\
    A: Managed identity where supported, otherwise Key Vault references or environment-injected secrets resolved at run time---never literals, because Git integration will eventually sync the notebook verbatim.
    \item \textbf{Q: Why batch sentiment calls instead of calling the API row-by-row?}\\
    A: Per-call latency and per-request pricing: batching amortizes both across rows, and bounded concurrency respects rate limits without serializing the cluster.
    \item \textbf{Q: What can Copilot in Fabric generate from natural language?}\\
    A: First-draft code and explanations---PySpark, KQL, API usage---which an engineer then reviews against the same correctness, leakage, and security standards as human code.
    \item \textbf{Q: Where should sentiment results land to be usable by BI?}\\
    A: In a Gold Delta table with a schema contract and \texttt{model\_version}, so Direct Lake models and Reflex rules consume them without recomputation.
    \item \textbf{Q: Why key the scoring delta on (\texttt{review\_id}, \texttt{model\_version})?}\\
    A: New reviews need scoring, and upgraded models legitimately re-score history; the composite key makes both behaviors correct and cheap.
\end{enumerate}

\section{Further Reading}
Review the SynapseML repository and documentation \cite{synapseML} and Azure OpenAI/AI services in Fabric \cite{microsoftFabricAiServices}. Copilot capabilities are documented at \cite{microsoftFabricCopilot}. The pipeline disciplines reused here come from Chapters~9--12 \cite{microsoftFabricPipelines,microsoftFabricMonitoringHub}; the Gold consumption path from Chapter~8 \cite{microsoftDirectLake}.

\section*{Key Takeaways}
\begin{itemize}
    \item SynapseML makes AI-service enrichment a distributed DataFrame operation, which is what terabyte-scale requires.
    \item Secrets come from Key Vault, identity, or environment---never cell text; Git integration makes ``temporary'' keys permanent.
    \item Batch, bound concurrency, and persist raw responses: the three plumbing rules of pipeline AI.
    \item Score deltas keyed on (entity, model version); re-scoring history is a deliberate, budgeted event.
    \item Copilot drafts, engineers own---review generated code for correctness, leakage, and secrets.
    \item Sentiment lands in Gold with a schema contract so BI and alerts consume without recompute.
    \item AI calls are metered dependencies: budget them, measure cost per thousand rows, and trend it.
\end{itemize}

\section{Exercises}
\begin{enumerate}
    \item \textbf{(Conceptual)} Compare the pre-built task APIs with generative endpoints for sentiment: cost, control, and fit.
    \item \textbf{(Conceptual)} Why does Git integration change the severity of a hardcoded key from bad habit to incident?
    \item \textbf{(Hands-on)} Extend the lab to a second language of reviews and document any label-distribution shift.
    \item \textbf{(Hands-on)} Implement the delta join on (\texttt{review\_id}, \texttt{model\_version}) and prove a rerun scores nothing.
    \item \textbf{(Design)} Design the retry-and-quarantine policy for the enrichment job: which failures retry, which quarantine, which page.
    \item \textbf{(Design)} Write the Copilot usage policy for the team: what it may draft, what review it must pass, what it may never touch.
    \item \textbf{(Interview)} In five sentences, explain the cost model of LLM enrichment to a product manager requesting ``real-time sentiment on everything.''
    \item \textbf{(Operational)} Write the runbook for a sudden spike in throttling errors from the AI service.
\end{enumerate}

\section{Capstone Project: Budgeted AI Enrichment Pipeline}\label{sec:ch19-capstone}
\index{capstone project}\index{SynapseML!capstone}\index{sentiment analysis!capstone}

\begin{capstonebox}
\textbf{Mission.} Productionize the Week~19 lab: a scheduled, delta-scoring sentiment pipeline over review data with vaulted secrets, bounded batches, quarantined failures, model-versioned Gold output, a cost report per run, and a Reflex alert on sustained negative-share climbs.

\textbf{Estimated time.} 5--7 hours.

\textbf{What you'll practice.}
\begin{itemize}
    \item Delta-keyed scoring with model-version awareness.
    \item Secret resolution and cost instrumentation in a scheduled pipeline.
    \item Failure quarantine and retry policy for a metered external dependency.
    \item Closing the loop from enrichment to BI and alerting.
\end{itemize}
\end{capstonebox}

\subsection*{Scenario}
Contoso Reviews aggregates product feedback. Engineering owns the sentiment pipeline end to end: nightly delta scoring, monthly cost reporting, and an alert when any category's negative share climbs for three consecutive days.

\subsection*{Requirements}
\begin{itemize}
    \item Scheduled pipeline scoring only the (\texttt{review\_id}, \texttt{model\_version}) delta.
    \item Secrets resolved at run time; zero literals in code or config.
    \item Quarantine table for persistently failing rows with error context.
    \item \texttt{gold.review\_sentiment} honoring the schema contract.
    \item A cost report: calls, estimated spend, cost per thousand rows, trend.
    \item The Reflex rule with sustained condition and cooldown, proven with injected negativity.
\end{itemize}

\subsection*{Implementation Steps}
\begin{enumerate}
    \item Refactor the lab notebook into a parameterized, secret-resolving scoring notebook.
    \item Implement the delta join and quarantine path.
    \item Schedule via pipeline with Chapter~12 alerting.
    \item Instrument and emit the cost report.
    \item Configure and prove the Reflex rule.
\end{enumerate}

\subsection*{Deliverables}
\begin{itemize}
    \item Scoring notebook and pipeline export in the repository.
    \item Delta-rerun evidence: second run scores zero rows.
    \item Quarantine evidence from an injected failure.
    \item Cost report for a week of runs.
    \item Reflex evidence from the injected negativity scenario.
\end{itemize}

\subsection*{Acceptance Criteria}
\begin{itemize}
    \item[\(\square\)] A rerun of the same batch scores zero rows; a model upgrade re-scores deliberately.
    \item[\(\square\)] No secret appears in any artifact, log, or screenshot.
    \item[\(\square\)] Persistent failures quarantine with context; transient failures retry with backoff.
    \item[\(\square\)] Gold output carries \texttt{model\_version} and passes the schema contract.
    \item[\(\square\)] The cost report computes per-thousand-row spend from measured calls.
    \item[\(\square\)] The Reflex alert fires once on the sustained scenario and not on noise.
\end{itemize}

\subsection*{Stretch Goals}
\begin{itemize}
    \item Add a second enrichment (key phrases) sharing the batch infrastructure.
    \item Compare task-API versus generative-endpoint cost at one million rows; publish the break-even.
    \item Add a daily Direct Lake sentiment dashboard consuming Gold.
    \item Add an evaluation harness: a labeled sample re-scored per model version to catch quality regressions.
\end{itemize}
